\section{CONCLUSION}

Bài báo này đã đề xuất thành công phương pháp mạng từ trở sinh lưới trong không gian ba chiều (3D-MBGRN) dựa trên hệ tọa độ trụ cho bài toán phân tích từ trường phi tuyến của máy điện quay đối xứng, với mô hình kiểm chứng là động cơ nam châm vĩnh cửu hướng trục (AFPM) 30 rãnh / 20 cực. Bằng việc rời rạc hóa miền tính toán thành tập hợp các phần tử hình lục diện đặc thù và áp dụng kỹ thuật dịch chỉ số lưới (index-shifting technique), phương pháp đề xuất đã mô hình hóa chuyển động quay của rotor một cách tự nhiên mà không đòi hỏi các tác vụ tái tạo lưới hay nội suy phức tạp tại mặt trượt. Bên cạnh đó, bộ giải phi tuyến lặp điểm cố định kết hợp thuật toán suy giảm thích nghi hệ số thư giãn (Adaptive Material Decay) và cơ chế quay lui nghiệm (rollback) đã giải quyết triệt để hiện tượng bẫy trì trệ và dao động số trị, đảm bảo tính hội tụ bền vững.

Kết quả kiểm chứng thực nghiệm đối chiếu trực tiếp với mô hình phần tử hữu hạn 3D (3D-FEM) khẳng định độ chính xác cao của bộ giải đề xuất trên cả hai trạng thái vận hành không tải và có tải định mức. Các đại lượng điện từ cốt lõi như mật độ từ thông khe hở ($B_z$), từ thông liên kết ($\Psi_a$), sức phản điện động ($e_a$), mô-men điện từ ($T$) và lực hút dọc trục ($F_z$) đều ghi nhận sự tương đồng rõ rệt về dạng sóng transient và phổ hài, với sai số đối chiếu chỉ giới hạn trong khoảng từ $0.3\%$ đến $3.1\%$.

Phân tích hội tụ lưới NRMSE và đánh giá tài nguyên tính toán cho thấy tại cấu hình lưới tối ưu Mesh 6 ($165.600$ ô lưới / $165.599$ DoF), bộ giải 3D-MBGRN đạt độ chính xác tiệm cận nghiệm chuẩn Mesh 9 (NRMSE $< 1.8\%$), đồng thời thời gian thực thi transient chỉ mất $0.3\text{ giờ}$ ($18\text{ phút}$), nhanh hơn $7.67\text{ lần}$ so với 3D-FEM ở cùng cấp độ Mesh 6 ($2.3\text{ giờ}$) và $18\text{ lần}$ so với 3D-FEM ở cấp độ mịn nhất Mesh 9 ($5.4\text{ giờ}$). Khả năng tiết kiệm tài nguyên tính toán và kiểm soát bộ nhớ RAM vượt trội này chứng minh bộ giải 3D-MBGRN hoàn toàn đủ năng lực thay thế các phần mềm 3D-FEM thương mại nặng nề trong các vòng lặp thiết kế tối ưu hóa đa tham số quy mô lớn.

Trong các nghiên cứu tiếp theo, bộ giải 3D-MBGRN sẽ tiếp tục được mở rộng để ứng dụng cho các cấu trúc máy điện 3D đối xứng phức tạp hơn (như máy điện nam châm chìm IPMSM hay máy điện dòng từ hỗn hợp Transverse Flux), đồng thời tích hợp trực tiếp bộ giải vào thuật toán tối ưu hóa đa mục tiêu (MOEA) nhằm đánh giá hiệu năng thực tế trong quy trình thiết kế động cơ tự động.